\chapter{Conclusions and outlook}
The central theme of this thesis has been the concrete tensor network realisation and study of various (generalised) symmetry structures. In conclusion, tensor networks provide a natural and effective framework in which these generalised symmetries can be explicitly represented on the lattice and manipulated. Not only do tensor networks provide a graphical calculus for dealing with these symmetries analytically, they also make it possible to put models with such symmetries on a computer and obtain their low-lying spectrum using variational methods. We reviewed this formalism from the ground up and illustrated it through several applications and research results.

Concretely, we saw that by phrasing topological order in the language of tensor networks, one can reveal hidden symmetry structures acting on the virtual degrees of freedom encoding the entanglement of the ground states. These symmetry structures are naturally described by fusion categories and realised as matrix product operators. Their relevance for the description of topological phases of matter can hardly be overstated. We demonstrated that they can be used naturally to construct distinct ground states as well as excited states, and argued that these symmetries persist beyond renormalisation fixed points.

In essence, these fusion categories describe simple symmetry operators, their fusion rules, and their anomalies. They therefore naturally generalise finite group and representation theory. Constructing all distinct representations of these categorical symmetries however requires bimodule categories. Although these may appear abstract at first sight, we have shown their importance in the construction of gapped boundaries of string-net models, and we have used them in the construction of a concrete quantum circuit that maps between Morita equivalent string-net models belonging to the same topological phase.

We then argued that these categorical symmetries also arise as internal symmetries of general Hamiltonians or classical models in one lower dimension. This holographic perspective allowed us to construct symmetric Hamiltonians, understand duality transformations as changes of module category, and relate the resulting dual models via bimodule tubes.

Finally, we showed that the same holographic tensor network perspective extends naturally to one higher dimension. We have shown that the symmetries of the (2+1)-dimensional models considered here are encoded in fusion 2-categories. This provides a higher-dimensional analogue of the categorical symmetry and duality framework developed earlier in the thesis.

\bigskip\noindent Within this framework, a number of open questions persist. Some of which, that spark my own interest, are the following.

One low-hanging fruit is to incorporate open boundary conditions into the symmetry and duality framework. In essence, this amounts to classifying the different ways in which symmetry operators can terminate at a physical boundary, and at the same time constructing the corresponding symmetric boundary Hamiltonian terms. This is particularly relevant for critical lattice models, where the infinite correlation length makes the bulk physics highly sensitive to the choice of boundary conditions~\cite{Fuchs2002,Thorngren2019,Choi2023,Lootens2022}.

From a technical point of view, module categories are expected to classify such boundary conditions in one dimension, and module 2-categories in two dimensions. However, it remains to understand how this classification interacts with the bulk degrees of freedom, and how to explicitly construct the corresponding boundary tensors for both the MPO and PEPO symmetries and the Hamiltonians.

Once such boundary conditions are incorporated, the next question is how to classify and explicitly construct the corresponding topological sectors, and how these sectors are mapped under duality transformations. This is particularly relevant for the numerical study of these open lattice models. Although the categorical structures expected to classify these sectors are formally known, turning this classification into an explicit tensor network construction remains an important open problem~\cite{Bullivant2020,Bullivant2021,Lootens2022}.

Also, it remains largely unexplored how to construct integrable lattice models with categorical symmetries within the tensor network formalism. In the main text, we only briefly touched upon the more natural setting for this question, namely two-dimensional classical lattice models. Such models can be constructed using the strange correlator construction which provides a systematic way of generating MPO symmetric transfer matrices depending on only a small number of parameters. The natural question is then which models within this low-dimensional family are (Yang-Baxter) integrable, and hence exactly solvable.

One promising approach is that of discrete holomorphicity, which imposes consistency equations on the Boltzmann weights guaranteeing the existence of a conserved current~\cite{Fendley2020}. It was shown how this naturally leads to known families of solutions to the Yang-Baxter equation. An interesting extension of the existing framework would be to allow currents described not by the input category itself, but by its Drinfeld center, thereby also covering input categories that are not braided. This is reminiscent of recent progress on the construction of critical lattice models with exotic symmetries~\cite{Huang2021,Vanhove2021b}.

By employing the tensor network formalism, it becomes possible to study these models numerically and potentially reconstruct the full conformal field theory spectrum from the obtained lattice models.

